\documentclass[10pt,a4paper]{amsart}
\usepackage[margin=19mm]{geometry}
\usepackage{amsmath,amssymb,mathtools}
\usepackage[scr=rsfs,cal=euler]{mathalfa}
\usepackage{xfrac}
\usepackage[unicode,hidelinks,bookmarks=false]{hyperref}
\newcommand{\ie}{i.e.\ }
\newcommand{\cf}{cf.\ }
\newcommand{\resp}{respectively\ }
\newcommand{\Subsection}{\subsection}
\providecommand{\nobreakdash}{\nobreak\hbox{-}}
\setlength{\emergencystretch}{2em}
\multlinegap0pt
\usepackage{amssymb}
\newtheorem{definition}{Definition}
\newtheorem{proposition}{Proposition}
\newtheorem{lemma}{Lemma}
\newtheorem{theorem}{Theorem}
\title{Ricci curvature and metric in causal spacetimes.}
\author{Javier Lafuente L\'{o}pez}
\date{Source: arXiv:2603.02027v1 (CC BY 4.0)}
\begin{document}
\maketitle

\noindent\emph{Source and licence.} Ricci curvature and metric in causal spacetimes.. Version arXiv:2603.02027v1 (CC BY 4.0), distributed by arXiv under the Creative Commons Attribution 4.0 International licence, http://creativecommons.org/licenses/by/4.0/, as recorded in the arXiv licence metadata for this exact version and shown on the abstract page https://arxiv.org/abs/2603.02027v1. Full licence notice: \texttt{licenses/arxiv-2603.02027-CC-BY-4.0.txt}.\par\smallskip

\noindent\textit{Editorial scope.} Counted unit: the article's theorem 1.8 (source0.tex lines 649--654). The article's complete original proof of it is delivered with it (source0.tex lines 656--703, one \texttt{proof} environment of 48 lines). No mathematics is written, repaired or re-derived: the statement and the proof are the article's own lines, in the article's own notation, and no retained line is substituted or re-flowed.  Retained uncounted as the source-local context the proof needs: lines 160--165; lines 282--290; lines 337--340; lines 454--458; lines 608--613.  Nothing was omitted from any retained span, so the package carries no omission record.  No retained line contains a citation, so the wrapper carries no bibliography and no bibliography entry.  No retained line contains a figure environment, an image-inclusion command or drawing code, so no image is delivered and the package declares no asset dependency.  Editorial formatting only, in addition: an AMS \texttt{amsart} preamble that restates the article's own declarations and macros used by the retained lines, namely the environments \texttt{definition}, \texttt{proposition}, \texttt{lemma}, \texttt{theorem} on separate counters, together with the standard packages those lines need.  The retained environments print in this document as Definition 1, Proposition 1, Lemma 1, Theorem 1 in order of appearance, whereas the article prints the same items with the numbering of its own full document.  The complete native source of the article as distributed in the arXiv e-print for this exact version is bundled unchanged as \texttt{sources/195590/source0.tex}.
\begin{definition}
\label{viable} The causal space $\left(  M,\mathcal{C}\right)  $ is said to be
viable, if there exists a metric $g\in\mathcal{C}$, such that $\left(
M,g\right)  $ is viable, in the sense that it has a timelike geodesic
$\gamma=\gamma\left(  t\right)  $ defined $\forall t>0$.
\end{definition}

\begin{proposition}
\label{0.5.3} The condition $\overline{\mathrm{Ric}}=\mathrm{Ric}$ is
equivalent to saying that the field $A$ satisfies%
\begin{equation}
\nabla_{X}A=\langle A,X\rangle A-\frac{1}{2}\langle A,A\rangle X,\quad\forall
X\in\mathfrak{X}(M)\label{ATP}%
\end{equation}

\end{proposition}

\begin{equation}
\nabla_{X}A=\langle A,X\rangle A-\frac{1}{2}\langle A,A\rangle X,\quad\forall
X\in\mathfrak{X}(M) \label{atp}%
\end{equation}

\begin{proposition}
\label{1.2}Let $A$ be an atypical field. If there exists $p\in M$ with
$\langle A,A\rangle(p)\neq0$, then it is $\langle A,A\rangle(x)\neq0$ for all
$x\in M$.
\end{proposition}

\begin{lemma}
\label{1.6}If $\overline{\mathrm{Ric}}=\mathrm{Ric}$, and $\langle
A,A\rangle=0$ on $M$. Then, if $A$ is not identically zero, then $A$ is
nowhere zero and $(M,g)$ is timelike and spacelike incomplete, and furthermore
all timelike $\nabla$-geodesics are incomplete.
\end{lemma}

\begin{theorem}
[Main.]\label{1.8}If $\left(  M,g\right)  $ is a viable spacetime (Definition
\ref{viable})and $\overline{\nabla}$ is a symmetric connection with the same
Ricci tensor as $g$ that preserves the light cones by parallel transport, then
$\overline{\nabla}$ is the metric connection of $g$.
\end{theorem}

\begin{proof}
If $\nabla$ is the metric connection of $g,$since $\overline{\mathrm{Ric}%
}=\mathrm{Ric}$ , by proposition \ref{0.5.3}, there exists $A$ atypic
vectorfield verifying (\ref{atp})

If $\langle A,A\rangle=0$, and there exists $p\in M$ with $A(p)\neq0$, by
Lemma \ref{1.6} $A$ does not vanish at any point, all timelike geodesics are
incomplete which contradicts the hypothesis. Suppose then that there exists
$p\in M$ with $\langle A,A\rangle(p)\neq0$. By \ref{1.2} we have $\langle
A,A\rangle\neq0$ in all of $M$. Let $\gamma:\mathbb{R}\longrightarrow M$ be
the timelike geodesic which by hypothesis is complete, and let $(E_{0}%
=\gamma^{\prime},E_{1},...,E_{n})$ be a parallel orthonormal basis along
$\gamma$ where $A=\sum a_{i}E_{i}$ with $a_{i}:\mathbb{R}\rightarrow
\mathbb{R}$ differentiable. Note that by Lemma 1.4 $A$ defines a field of
incomplete pregeodesics, so $A\circ\gamma$ is never tangent to $\gamma$, that
is, the function:
\[
f(t)=\frac{1}{2}(a_{1}(t)^{2}+...+a_{n}(t)^{2})
\]
verifies $f(t)>0$ for all $t\in\mathbb{R}$. Since $\frac{\nabla A}{dt}%
=\sum\frac{da_{i}}{dt}E_{i}=\langle A,\gamma^{\prime}\rangle A-\frac{1}%
{2}\langle A,A\rangle\gamma^{\prime}=-a_{0}[\sum a_{i}E_{i}]-\frac{1}%
{2}(-a_{0}^{2}+a_{1}^{2}+...+a_{n}^{2})E_{0}$, and comparing the coefficients
multiplying $E_{0}$ we get:
\[
\frac{da_{0}}{dt}=-\frac{1}{2}a_{0}^{2}-f(t)
\]


Calling $y(t)=-a_{0}(t)$, we conclude that $y(t)$ is defined on all of
$\mathbb{R}$ and satisfies the differential equation:
\begin{equation}
\frac{dy}{dt}=\frac{1}{2}y^{2}+f(t) \label{EDO}%
\end{equation}
where $f:\mathbb{R}\longrightarrow\mathbb{R}$ is differentiable, with
$f(t)>0$.
\[%
\begin{tabular}
[c]{l}%
We can choose the direction of traversal of\\
$\gamma$, such that $a_{0}\left(  0\right)  =-\left.  \left\langle
A,E_{0}\right\rangle \right\vert _{t=0}>0$%
\end{tabular}
\]


This contradicts the following Lemma.
\end{proof}

\end{document}
